\section{Algebra and numerics (KACTL)}

\entry{Discrete log, modular sqrt}{KACTL \quad $O(\sqrt m)$, $O(\log^2 p)$}{%
\code{modLog(a, b, m)}: the smallest $x > 0$ with $a^x \equiv b \pmod m$, or \code{-1}.
Any \code{m}, prime or not, \code{a} need not be coprime to it. Baby-step giant-step:
store $b\cdot a^j$ for $j \le \sqrt m$, then step by $a^{\sqrt m}$ and look each up.
$x = 0$ is never returned, so if $b \equiv 1$ is a valid answer at exponent 0, check that
yourself; \code{modLog(a, 1, m)} is the \emph{order} of \code{a}. Needs
$m < 3\cdot10^9$ (\code{e*a} must fit).\\[0.3ex]
\code{modSqrt(a, p)}: an $x$ with $x^2 \equiv a \pmod p$, \textbf{p prime}; the other
root is \code{p - x}. Returns \code{-1} if \code{a} is not a square (Euler's criterion,
$a^{(p-1)/2} \ne 1$) --- KACTL asserts instead. Tonelli--Shanks. Uses \code{power}
from \emph{Modular basics}.}

%DEP modbasics
%LABEL modlog
\begin{lstlisting}
ll modLog(ll a, ll b, ll m) {
    ll n = (ll) sqrt(m) + 1, e = 1, f = 1, j = 1;
    unordered_map<ll, ll> A;
    while (j <= n && (e = f = e * a % m) != b % m)
        A[e * b % m] = j++;
    if (e == b % m) return j;
    if (__gcd(m, e) == __gcd(m, b))
        rep(i,2,n+2) if (A.count(e = e * f % m))
            return n * i - A[e];
    return -1;
}
ll modSqrt(ll a, ll p) {          // KACTL: sqrt(a, p)
    a %= p; if (a < 0) a += p;
    if (a == 0) return 0;
    if (power(a, (p-1)/2, p) != 1) return -1;  // no root
    if (p % 4 == 3) return power(a, (p+1)/4, p);
    ll s = p - 1, n = 2;
    int r = 0, m;
    while (s % 2 == 0)
        ++r, s /= 2;
    // find a non-square mod p
    while (power(n, (p - 1) / 2, p) != p - 1) ++n;
    ll x = power(a, (s + 1) / 2, p);
    ll b = power(a, s, p), g = power(n, s, p);
    for (;; r = m) {
        ll t = b;
        for (m = 0; m < r && t != 1; ++m)
            t = t * t % p;
        if (m == 0) return x;
        ll gs = power(g, 1LL << (r - m - 1), p);
        g = gs * gs % p;
        x = x * gs % p;
        b = b * g % p;
    }
}
\end{lstlisting}

\entry{FFT, convolution mod any m}{KACTL \quad $O(n \log n)$}{%
Same job as the NTT entry, \code{c[i] = sum a[j]*b[i-j]}, for the two cases NTT cannot
do.\\[0.3ex]
\code{conv(a, b)} on \code{double}s: \textbf{exact integer} results without any
modulus, e.g.\ counting pairs whose sum is $s$ when the counts reach $10^{10}$ (past
NTT's prime). Round with \code{llround}. Safe while
$(\sum a_i^2 + \sum b_i^2)\log_2 N < 9\cdot10^{14}$.\\[0.3ex]
\code{convMod<M>(a, b)}: the answer modulo \textbf{any} \code{M} up to $\sim10^9$ ---
the usual case is $10^9{+}7$, where NTT does not work because $10^9{+}6$ has no large
power of two. Inputs must already be in \code{[0, M)}. It splits every value into
$\sqrt M$-sized halves, so it is about twice as slow as one NTT.\\[0.3ex]
KACTL names the complex type \code{C}; it is \code{cd} here because \code{C(n, k)} is
the binomial.}

%LABEL fft
\begin{lstlisting}
typedef complex<double> cd;
typedef vector<double> vd;
void fft(vector<cd>& a) {
    int n = (int)a.size(), L = 31 - __builtin_clz(n);
    static vector<complex<long double>> R(2, 1);
    static vector<cd> rt(2, 1);  // (^ 10% faster if double)
    for (static int k = 2; k < n; k *= 2) {
        R.resize(n); rt.resize(n);
        auto x = polar(1.0L, acos(-1.0L) / k);
        rep(i,k,2*k) rt[i] = R[i] = i&1 ? R[i/2] * x : R[i/2];
    }
    vi rev(n);
    rep(i,0,n) rev[i] = (rev[i / 2] | (i & 1) << L) / 2;
    rep(i,0,n) if (i < rev[i]) swap(a[i], a[rev[i]]);
    for (int k = 1; k < n; k *= 2)
        for (int i = 0; i < n; i += 2 * k) rep(j,0,k) {
            cd z = rt[j+k] * a[i+j+k];
            a[i + j + k] = a[i + j] - z;
            a[i + j] += z;
        }
}
vd conv(const vd& a, const vd& b) {
    if (a.empty() || b.empty()) return {};
    vd res(a.size() + b.size() - 1);
    int L = 32 - __builtin_clz((int)res.size()), n = 1 << L;
    vector<cd> in(n), out(n);
    copy(all(a), begin(in));
    rep(i,0,(ll)b.size()) in[i].imag(b[i]);
    fft(in);
    for (cd& x : in) x *= x;
    rep(i,0,n) out[i] = in[-i & (n - 1)] - conj(in[i]);
    fft(out);
    rep(i,0,(ll)res.size()) res[i] = imag(out[i]) / (4 * n);
    return res;
}
template<int M> vi convMod(const vi &a, const vi &b) {
    if (a.empty() || b.empty()) return {};
    vi res(a.size() + b.size() - 1);
    int B = 32 - __builtin_clz((int)res.size()), n = 1 << B;
    int cut = int(sqrt(M));
    vector<cd> L(n), R(n), outs(n), outl(n);
    rep(i,0,(ll)a.size())
        L[i] = cd((int)a[i] / cut, (int)a[i] % cut);
    rep(i,0,(ll)b.size())
        R[i] = cd((int)b[i] / cut, (int)b[i] % cut);
    fft(L), fft(R);
    rep(i,0,n) {
        int j = -i & (n - 1);
        outl[j] = (L[i] + conj(L[j])) * R[i] / (2.0 * n);
        outs[j] = (L[i] - conj(L[j])) * R[i] / (2.0 * n) / 1i;
    }
    fft(outl), fft(outs);
    rep(i,0,(ll)res.size()) {
        ll av = ll(real(outl[i])+.5);
        ll cv = ll(imag(outs[i])+.5);
        ll bv = ll(imag(outl[i])+.5) + ll(real(outs[i])+.5);
        res[i] = ((av % M * cut + bv) % M * cut + cv) % M;
    }
    return res;
}
\end{lstlisting}

\entry{AND / OR / XOR convolution}{KACTL \quad $O(n \log n)$}{%
\code{c[z] = sum over x op y == z of a[x]*b[y]}, with op one of AND, OR, XOR --- the
bitwise cousin of FFT. Typical: ``how many pairs have XOR exactly $k$'' (put the counts
in \code{a} and \code{b}, read \code{c[k]}), or combining two independent choices
described by bitmasks. Pad both to the same power of two.\\[0.3ex]
\textbf{Pick the op by leaving exactly one of the three \code{tie} lines active}; for
XOR also uncomment the final division. Mod a prime: reduce after every \code{+}/\code{-},
and for XOR multiply by $n^{-1}$ instead of dividing. Called \code{conv} in KACTL,
renamed so it does not clash with NTT's.}

%LABEL fst
\begin{lstlisting}
void FST(vi& a, bool inv) {
    for (int n = (int)a.size(), step = 1; step < n; step *= 2)
        for (int i = 0; i < n; i += 2*step) rep(j,i,i+step) {
            ll &u = a[j], &v = a[j + step]; tie(u, v) =
                inv ? pii(v - u, u) : pii(v, u + v); // AND
             // inv ? pii(v, u - v) : pii(u + v, u); // OR
             // pii(u + v, u - v);                   // XOR
        }
    // if (inv) for (ll& x : a) x /= (ll)a.size(); // XOR only
}
vi fstConv(vi a, vi b) {
    FST(a, 0); FST(b, 0);
    rep(i,0,(ll)a.size()) a[i] *= b[i];
    FST(a, 1); return a;
}
\end{lstlisting}

\entry{Berlekamp--Massey + k-th term}{KACTL \quad $O(n^2)$, $O(n^2 \log k)$}{%
The ``I can compute the first 50 answers but need the $10^{18}$-th'' tool. If a sequence
mod a prime satisfies \emph{some} linear recurrence
$S_i = \sum_j tr_j\, S_{i-j-1}$ of order $\le n$, then \code{berlekampMassey} recovers
\code{tr} from the first $2n$ terms. Any DP whose step is a fixed linear map on a
fixed-size state has one (tilings of a $k\times n$ board, walks in a fixed graph,
counting strings avoiding a pattern), of order at most the state size, so you never derive
the matrix: brute-force the first terms, feed them in.\\[0.3ex]
Then \code{linearRec(S, tr, k)} gives term $k$ (0-indexed) from the first
\code{tr.size()} terms --- faster than matrix power. Example:
\code{berlekampMassey(\{0,1,1,2,3,5\})} $=$ \code{\{1,1\}}, and
\code{linearRec(\{0,1\}, \{1,1\}, k)} is Fibonacci.\\[0.3ex]
Give it \textbf{plenty} of terms: if the result's length is close to half of what you
supplied, the recurrence may be longer --- add terms until it stops changing. An all-zero
input gives an empty \code{tr}; handle that before calling \code{linearRec}. Uses
\code{power} from \emph{Modular basics}; \code{mod} is shared with the next entry.}

%DEP modbasics
%LABEL bm
\begin{lstlisting}
const ll mod = 1000000007;          // must be prime for BM
vi berlekampMassey(vi s) {
    int n = (int)s.size(), L = 0, m = 0;
    vi C(n), B(n), T;
    C[0] = B[0] = 1;
    ll b = 1;
    rep(i,0,n) { ++m;
        ll d = s[i] % mod;
        rep(j,1,L+1) d = (d + C[j] * s[i - j]) % mod;
        if (!d) continue;
        T = C; ll coef = d * power(b, mod-2, mod) % mod;
        rep(j,m,n) C[j] = (C[j] - coef * B[j - m]) % mod;
        if (2 * L > i) continue;
        L = i + 1 - L; B = T; b = d; m = 0;
    }
    C.resize(L + 1); C.erase(C.begin());
    for (ll& x : C) x = (mod - x) % mod;
    return C;
}
ll linearRec(vi S, vi tr, ll k) {
    int n = (int)tr.size();
    auto combine = [&](vi a, vi b) {
        vi res(n * 2 + 1);
        rep(i,0,n+1) rep(j,0,n+1)
            res[i + j] = (res[i + j] + a[i] * b[j]) % mod;
        for (int i = 2 * n; i > n; --i) rep(j,0,n)
            res[i-1-j] = (res[i-1-j] + res[i] * tr[j]) % mod;
        res.resize(n + 1);
        return res;
    };
    vi pol(n + 1), e(pol);
    pol[0] = e[1] = 1;
    for (++k; k; k /= 2) {
        if (k % 2) pol = combine(pol, e);
        e = combine(e, e);
    }
    ll res = 0;
    rep(i,0,n) res = (res + pol[i + 1] * S[i]) % mod;
    return res;
}
\end{lstlisting}

\entry{Determinant}{KACTL \quad $O(n^3)$}{%
Both versions \textbf{destroy} the matrix --- pass a copy if you need it again.
\code{double}: Gaussian elimination with partial pivoting. \code{ll}: the exact answer
mod \code{mod}, and \code{mod} need \emph{not} be prime, because it eliminates with
Euclid-style row subtractions instead of inverses (remove every \code{\% mod} for plain
integers if nothing overflows).\\[0.3ex]
Main use, \textbf{Kirchhoff's matrix-tree theorem}: the number of spanning trees is the
determinant of the Laplacian (\code{L[v][v]} = degree, \code{L[u][v]} = $-$(number of
edges $uv$)) with any one row and the same column deleted. Also the
Lindstr\"om--Gessel--Viennot lemma: the number of vertex-disjoint path systems in a DAG
is the determinant of the path-count matrix.}

%DEP modbasics bm
%LABEL det
\begin{lstlisting}
double det(vector<vector<double>>& a) {
    int n = (int)a.size(); double res = 1;
    rep(i,0,n) {
        int b = i;
        rep(j,i+1,n) if (fabs(a[j][i]) > fabs(a[b][i])) b = j;
        if (i != b) swap(a[i], a[b]), res *= -1;
        res *= a[i][i];
        if (res == 0) return 0;
        rep(j,i+1,n) {
            double v = a[j][i] / a[i][i];
            if (v != 0) rep(k,i+1,n) a[j][k] -= v * a[i][k];
        }
    }
    return res;
}
ll det(vector<vector<ll>>& a) {          // mod `mod`
    int n = (int)a.size(); ll ans = 1;
    rep(i,0,n) {
        rep(j,i+1,n) {
            while (a[j][i] != 0) { // gcd step
                ll t = a[i][i] / a[j][i];
                if (t) rep(k,i,n)
                    a[i][k] = (a[i][k] - a[j][k] * t) % mod;
                swap(a[i], a[j]);
                ans *= -1;
            }
        }
        ans = ans * a[i][i] % mod;
        if (!ans) return 0;
    }
    return (ans + mod) % mod;
}
\end{lstlisting}

\entry{Simplex (linear programming)}{KACTL \quad fast in practice, $2^n$ worst}{%
Maximise $c^T x$ subject to $Ax \le b$, $x \ge 0$ (real $x$). Returns the optimum and
sets \code{x}; \code{-inf} if infeasible, \code{inf} if unbounded. Recognise it when the
problem is ``choose real amounts, linear limits, linear objective'' (diets, mixing,
fractional scheduling) and $n, m \lesssim$ a few hundred.\\[0.3ex]
Rewriting into that form: \textbf{minimise} $\to$ negate \code{c} and the result;
$\ge$ row $\to$ negate the row and its \code{b}; \textbf{equality} $\to$ both $\le$ and
$\ge$; a \textbf{free} variable $\to$ $x = x^+ - x^-$, two columns. The answer is a
\code{double}; it is not an integer program, rounding \code{x} is not valid in
general.\\[0.3ex]
Usage: \code{vvd A = \{\{1,-1\},\{-1,1\},\{-1,-2\}\}; vd b = \{1,1,-4\}, c = \{-1,-1\}, x;}
\code{T val = LPSolver(A, b, c).solve(x);}\\[0.3ex]
\code{eps} renamed to \code{EPS} because the header defines \code{eps}.}

%LABEL simplex
\begin{lstlisting}
typedef double T; // long double, Rational, double + mod<P>
typedef vector<T> vd;
typedef vector<vd> vvd;
const T EPS = 1e-8, inf = 1/.0;
#define MP make_pair
#define ltj(X) if(s==-1 || MP(X[j],N[j]) < MP(X[s],N[s])) s=j
struct LPSolver {
    int m, n;
    vi N, B;
    vvd D;
    LPSolver(const vvd& A, const vd& b, const vd& c) :
        m((int)b.size()), n((int)c.size()),
        N(n+1), B(m), D(m+2, vd(n+2)) {
        rep(i,0,m) rep(j,0,n) D[i][j] = A[i][j];
        rep(i,0,m) { B[i] = n+i; D[i][n] = -1;
                     D[i][n+1] = b[i]; }
        rep(j,0,n) { N[j] = j; D[m][j] = -c[j]; }
        N[n] = -1; D[m+1][n] = 1;
    }
    void pivot(int r, int s) {
        T *a = D[r].data(), inv = 1 / a[s];
        rep(i,0,m+2) if (i != r && abs(D[i][s]) > EPS) {
            T *b = D[i].data(), inv2 = b[s] * inv;
            rep(j,0,n+2) b[j] -= a[j] * inv2;
            b[s] = a[s] * inv2;
        }
        rep(j,0,n+2) if (j != s) D[r][j] *= inv;
        rep(i,0,m+2) if (i != r) D[i][s] *= -inv;
        D[r][s] = inv;
        swap(B[r], N[s]);
    }
    bool simplex(int phase) {
        int x = m + phase - 1;
        for (;;) {
            int s = -1;
            rep(j,0,n+1) if (N[j] != -phase) ltj(D[x]);
            if (D[x][s] >= -EPS) return true;
            int r = -1;
            rep(i,0,m) {
                if (D[i][s] <= EPS) continue;
                if (r == -1 || MP(D[i][n+1] / D[i][s], B[i])
                            < MP(D[r][n+1] / D[r][s], B[r]))
                    r = i;
            }
            if (r == -1) return false;
            pivot(r, s);
        }
    }
    T solve(vd &x) {
        int r = 0;
        rep(i,1,m) if (D[i][n+1] < D[r][n+1]) r = i;
        if (D[r][n+1] < -EPS) {
            pivot(r, n);
            if (!simplex(2) || D[m+1][n+1] < -EPS)
                return -inf;
            rep(i,0,m) if (B[i] == -1) {
                int s = 0;
                rep(j,1,n+1) ltj(D[i]);
                pivot(i, s);
            }
        }
        bool ok = simplex(1); x = vd(n);
        rep(i,0,m) if (B[i] < n) x[B[i]] = D[i][n+1];
        return ok ? D[m][n+1] : inf;
    }
};
\end{lstlisting}

\entry{Numerical integration}{KACTL \quad $O(n)$ evaluations}{%
Simpson's rule on $2n$ slices, error $\propto h^4$. For areas with no closed form:
union of circles (integrate the covered length of each vertical line), volumes of
revolution. Re-run with a bigger \code{n} and check the digits you need stay put; if
\code{f} has kinks (e.g.\ where circles start and end), integrate piecewise between
them instead of one big call.}

%LABEL quad
\begin{lstlisting}
template<class F>
double quad(double a, double b, F f, const int n = 1000) {
    double h = (b - a) / 2 / n, v = f(a) + f(b);
    rep(i,1,n*2)
        v += f(a + i*h) * (i&1 ? 4 : 2);
    return v * h / 3;
}
\end{lstlisting}
